\begin{textsample}{NEG Dim 9 – global\_south\_2024\_10 – Score -5.00 – global\_south\_2024\_10/115567118.txt}  \label{ex:f9_neg_280}
Philippine embassy pays tribute to victims of Hamas attack Prayers were offered yesterday to remember Filipino caregivers Loreta Alacre, Angelyn Aguirre, Grace Cabrera and Paul Vincent Castelvi who were \textbf{murdered} in the attack, which also led to the abduction of more than 250 persons.
Philippine Embassy in Israel / Facebook page MANILA, Philippines -- On the first anniversary of the Oct. 7 Hamas-led terror attack in Israel, the Philippine embassy in Tel Aviv and the Filipino community remembered and paid tribute to the 1,200 Israelis and foreign nationals, including four Filipinos, who were brutally killed.
Prayers were offered yesterday to remember Filipino caregivers Loreta Alacre, Angelyn Aguirre, Grace Cabrera and Paul Vincent Castelvi who were \textbf{murdered} in the attack, which also led to the abduction of more than 250 persons.
Alacre was killed while attending a \textbf{music} \textbf{festival} near the Gaza Strip.
Aguirre, Cabrera and Castelvi were \textbf{murdered} along with their patients and employers by Hamas terrorists in Beeri, a \textbf{kibbutz} next to Gaza.
Fleur Hassan-Nahoum, deputy mayor of Jerusalem, honored Aguirre @ @ @ @ @ @ @ @ @ @ for elderly patient Nira in \textbf{Kibbutz} Kfar Aza, near Gaza Strip."Despite having a chance to flee the Hamas terror attacks, ( Aguirre ) showed unbelievable humanity and loyalty by remaining ( at ) Nira ' s side during the violence, resulting in both of them being brutally \textbf{murdered} by Hamas,"Hassan-Nahoum said."May their memories be blessings forever,"the embassy said."We continue to call for the release of the remaining hostages, offer comfort for the bereaved families and advocate for a peaceful resolution to the ongoing conflict, toward lasting peace in the region,"it added.
Meanwhile, the embassy advised Filipinos in Israel to avoid the Temple Mount, Damascus Gate, Herod ' s Gate, Al Wad Road, Musrara Road and the vicinity of East Jerusalem due to security concerns amid the tension in the Middle East.
The advisory also included the West Bank ( areas near the border of Gaza and Lebanon ) and Golan Heights ( areas near sensitive installations such @ @ @ @ @ @ @ @ @ @ ).
Filipinos are also urged to avoid crowded places like malls and markets for non-essential activities and big gatherings, such as \textbf{concert}, rally and \textbf{festivals}."At the first sign of danger, evacuate immediately and go to a safe place.
If this place is nearby, determine if it is possible to go there by bicycle, electric scooter or by walking,"a part of the advisory reads.
Filipinos were also advised not to approach Israeli security forces in sensitive areas.

% matched lemmas: concert, festival, kibbutz, murder, music
\end{textsample}
